\section{Ground Penetrating Radar (GPR)}
\label{subsec:groud_penetrating_radar}

O método aqui descrito é fundamentado no envio de um pulso eletromagnético por meio de uma antena transmissora. Esse pulso tem uma duração extremamente curta de apenas alguns nanosegundos ($10^{-9}$ segundos) empregando ondas eletromagnéticas de frequências VHF/UHF (conforme ilustrado na figura \ref{fig:fig02}).

\begin{figure}[H]
  \centering
  \includegraphics[width=1\linewidth]{lib/fig/gpr/fig01.png}
  % Contém Fonte: 
  \caption{O espectro eletromagnético delineia as gamas de frequência utilizadas por várias tecnologias. O GPR opera dentro da faixa de frequência de 10 MHz a 1 GHz, correspondendo a comprimentos de onda que variam de aproximadamente 30 metros a 0.3 metros.}
\source{\cite{xavier2006processamento}}
  \label{fig:fig01}
\end{figure}

Assim como os métodos sísmicos, o método \ac*{gpr} opera com o registro do tempo que leva para um pulso eletromagnético viajar. Esse pulso é emitido a partir da fonte transmissora (TX), penetra no subsolo e encontra estruturas com contrastes de impedância eletromagnética suficientes para causar reflexões. Em seguida, o pulso retorna ao receptor (RX). A figura 2.2 ilustra o funcionamento do equipamento em um cenário com duas camadas de constantes dielétricas, K1 e K2. A figura também mostra o traço que representa a onda direta pelo ar (Aw), a onda direta através da terra (Gw) e o registro do tempo de viagem duplo até a interface.

\begin{figure}[H]
  \centering
  \includegraphics[width=1\linewidth]{lib/fig/gpr/fig02.png}
  % Contém Fonte: 
  \caption{(A) Princípio de funcionamento do método \ac*{gpr} considerando um ambiente
composto por duas camadas. (B) Traço contendo os registros dos tempos de chegada da
onda direta pelo ar (Aw), on da direta pela terra (Gw) e da onda refletida na interface
formada pelas camadas 1 e 2.}
\source{\cite{oliveira2001dois}}
  \label{fig:fig02}
\end{figure}

Quando esse pulso atinge diferentes meios (ou alvos) com distintas propriedades eletromagnéticas, a velocidade de propagação é alterada. Como resultado, parte da energia eletromagnética é refletida e captada por uma antena receptora, enquanto outra parte sofre difração e continua se propagando até encontrar outra superfície de contraste ou ser completamente atenuada pelo meio.

O "tempo duplo" é calculado, transmitido, registrado e armazenado por uma unidade de controle. Esse tempo duplo refere-se ao intervalo decorrido entre a emissão do pulso eletromagnético pela antena transmissora até o momento em que seu retorno é registrado pela antena receptora, após atingir o alvo e retornar à superfície. Além disso, a unidade de controle também registra a intensidade correspondente do sinal refletido.


\begin{figure}[H]
  \centering
  \includegraphics[width=1\linewidth]{lib/fig/gpr/fig03.png}
  % Contém Fonte: 
  \caption{Esquema de utilização do GPR com transmissão e recebimento de sinais.}
\source{\cite{sampaio2021investigaccao}}
  \label{fig:fig03}
\end{figure}

Dessa forma, é possível obter imagens que revelam informações sobre as condições geomecânicas do terreno ou superfície em estudo, como a densidade de fraturas existentes e seu grau de fraturamento. Além disso, é capaz de determinar o topo do maciço rochoso e identificar descontinuidades, espessura do pacote de sedimentos, textura e nível d'água \cite{davis1989ground}.
Na figura \ref{fig:fig03} mostra o funcionamento do \ac*{gpr}. \cite{annan1992ground}

A validação dos resultados obtidos com o radar evidencia o poder e utilidade desse método como uma ferramenta auxiliar na interpolação e extrapolação das informações obtidas por meio de métodos de perfuração convencionais. \cite{davis1989ground}

\subsection{Janela de GPR}
\label{subsubsec:janela_gpr}

A faixa de operação do \ac*{gpr} consiste em um intervalo de frequência no qual a atenuação experimentada pelo sinal de radar permanece constante para todos os materiais presentes no substrato geológico em consideração. É precisamente dentro dessa faixa que o fenômeno de propagação do pulso eletromagnético é governado pela equação de onda. A frequência associada ao limite inferior dessa faixa deve ser suficientemente alta para evitar a difusão da onda. Por outro lado, o limite superior não pode exceder 1 GHz, uma vez que isso levaria ao fenômeno de polarização das moléculas de água presentes nas rochas, resultando em uma significativa atenuação do sinal \cite{annan1992ground}.

\begin{figure}[H]
  \centering
  \includegraphics[width=1\linewidth]{lib/fig/gpr/fig04.png}
  % Contém Fonte: 
  \caption{A relação entre a velocidade e a frequência em diferentes condutividades.}
\source{\cite{annan1992ground}}
  \label{fig:fig04}
\end{figure}

As velocidades dos sinais de radar em materiais geológicos de baixa perda, adequados para sondagens por radar, estão relacionadas com a parte real da constante dielétrica pela fórmula:
\[V = \frac{c}{(K')^{1/2}}\],
onde \(c = 3 \times 10^8 \, \text{m/s}\), a velocidade de propagação de ondas eletromagnéticas no vácuo.

\begin{table}[H]
    \centering
    \caption{Constante dielétrica típica, condutividade elétrica, velocidade
e atenuação observadas em materiais geológicos comuns a 100
MHz.}

    \begin{tabular}{@{}ccccc@{}}
        \toprule
        \textbf{Material}  & \textbf{Constante
Dielétrica} & \textbf{Condutividade} & \textbf{Velocidade
}      & \textbf{Atenuação}      \\
        \textbf{}  & \textbf{K} & \textbf{$\sigma$(mS/m)} & \textbf{$\nu$(m/ns)
}      & \textbf{$\alpha$(dB/m)}      \\ \midrule
        Ar & 1 & 0 & 0.3 & 0 \\
        Água destilada & 80 & 0.01 & 0.033 & 2 x $10^{-3}$ \\
        Água fresca & 80 & 0.5 & 0.033 & 0.1 \\
        Água do mar & 80 & 3 x $10^{4}$ & 0.01 & 103 \\
        Areia fresca & 3-5 & 0.01 & 0.15 & 0.01 \\
        Areia saturada & 20-30 & 0.1-1 & 0.06 & 0.03-0.3 \\
        Calcário & 4-8 & 0.5-2 & 0.12 & 0.4-1 \\
        Folhelho & 5-15 & 1-100 & 0.09 & 1-100 \\
        Siltito & 5-30 & 1-100 & 0.07 & 1-100 \\
        Argila & 5-40 & 2-1000 & 0.06 & 1-300 \\
        Granito & 4-6 & 0.01-1 & 0.13 & 0.01-1 \\
        Sal & 5-6 & 0.01-1 & 0.13 & 0.01-1 \\
        Gelo & 3-4 & 0.01 & 0.16 & 0.01 \\
        \multicolumn{1}{l}{} & \multicolumn{1}{l}{}         & \multicolumn{1}{l}{} & \multicolumn{1}{l}{} \\ \bottomrule
    \end{tabular}
    \source{\cite{daniels2004ground}}
    \label{tab:tab1}
\end{table}

\subsection{GPRPy}
\label{subsubsec:gprpy}

Muitas vezes, as investigações envolvendo \ac{gpr} não exigem um processamento avançado para obter resultados úteis. No entanto, a maioria dos estudos de pesquisa em \ac{gpr} utiliza software proprietário fechado fornecido pelos fabricantes de sistemas \ac{gpr} ou pacotes independentes \citeonline{jol2008ground}. Essa dependência em software proprietário apresenta várias desvantagens significativas. Em primeiro lugar, o software oferece pouca flexibilidade, impedindo os usuários de adaptá-lo completamente às suas necessidades específicas. Em segundo lugar, a capacidade de reproduzir pesquisas é limitada, pois para replicar os passos de processamento mencionados em um artigo, os pesquisadores precisam ter acesso ao mesmo software proprietário. Essa limitação se torna mais crítica à medida que editores e agências de financiamento pressionam por pesquisas transparentes. Em terceiro lugar, a disponibilidade de cursos que ensinam o processamento de \ac{gpr} também é prejudicada devido à dependência do software proprietário.

O pacote de software \ac{gprpy} consiste em dois interfaces gráficas de usuário (GUIs): uma para o processamento de dados de perfil figura \ref{fig:fig07} e outra para dados de ponto médio comum (CMP) ou reflexão e refração de amplo ângulo (WARR) figura \ref{fig:fig07}. O pacote também inclui uma série de funções para processamento mais avançado, como a geração de cubos de dados 3D por interpolação de perfis, interpolação de superfícies entre pontos selecionados e fusão de perfis. O código de computador da GUI é separado do código de computador que implementa o processamento e a visualização de dados. Essa separação permite que os usuários escrevam código Python que acesse diretamente todas as capacidades do \ac{gprpy} sem precisar usar a GUI. O \ac{gprpy} aproveita essa separação por meio de seus scripts gerados automaticamente.

\begin{figure}[H]
  \centering
  \includegraphics[width=1\linewidth]{lib/fig/gpr/fig07.png}
  % Contém Fonte: 
  \caption{Interface Gráfica do Usuário (GUI) para o processamento de dados de perfil de ponto médio comum (CMP) no \ac{gpr}.}
\source{\cite{plattner2020gprpy}}
  \label{fig:fig07}
\end{figure}

\begin{figure}[H]
  \centering
  \includegraphics[width=1\linewidth]{lib/fig/gpr/fig08.png}
  % Contém Fonte: 
  \caption{Interface Gráfica do Usuário (GUI) para a análise de velocidade a partir de dados CMP ou WARR, com representações de empilhamento de amplitudes para padrões de tempo de chegada em forma de hipérbole (a) e em forma linear (b), bem como os próprios dados CMP ou WARR (c).}
\source{\cite{plattner2020gprpy}}
  \label{fig:fig08}
\end{figure}

\section{ANÁLISE E INTERPRETAÇÃO DE DADOS}
\label{subsubsec:analiseeinterpretação}

O processamento de dados do \ac*{gpr} pode seguir dois caminhos principais. O primeiro envolve a apresentação dos dados em forma de seções, planos ou volumes para indicar a localização de alvos anômalos, comum a todos os métodos geofísicos. O segundo caminho consiste em extrair variáveis das propriedades das ondas, como velocidade, atenuação ou impedância, e convertê-las em informações específicas para a aplicação em questão.

O processo típico de processamento de dados do \ac*{gpr} inclui etapas de edição de dados, processamento básico, processamento avançado e visualização/interpretação. O processamento é uma atividade iterativa, e os dados podem passar várias vezes por esse ciclo. Em grandes conjuntos de dados, pode ser aplicado um processamento em lote com controle interativo limitado, após testes iniciais em amostras selecionadas.

\begin{figure}[H]
  \centering
  \includegraphics[width=1\linewidth]{lib/fig/gpr/fig05.png}
  % Contém Fonte: 
  \caption{Visão geral do fluxo de processamento de dados do radar de penetração no solo (GPR).}
\source{\cite{jol2008ground}}
  \label{fig:fig05}
\end{figure}

\begin{itemize}
    \item Instrumento: Equipamentos ou dispositivos usados para coletar dados ou medições;
    \item Impressão: Gravação de dados que são exibidos ou impressos de alguma forma;
    \item Exibição/Impressão: Exibir ou imprimir dados, para fins de visualização ou análise;
    \item Visualização: Ferramentas ou software usados para criar representações visuais de dados, tornando mais fácil compreendê-los e analisá-los;
    \item Editar: Processo de edição ou modificação de dados, o que pode ser necessário para limpar ou pré-processar os dados;
    \item Processo Básico: Etapas fundamentais de processamento de dados, como filtragem, redução de ruído ou cálculos simples;
    \item Processo Avançado: Técnicas de processamento de dados mais complexas que vão além das operações básicas. Pode incluir tarefas como extração de características, reconhecimento de padrões ou análise estatística;
    \item Processamento de Imagens: Imagem processada, após tratamento básico e avançado, gerando imagem imagem geológica ou não, para uma melhor interpretação dos dados;
    
\end{itemize}

\begin{algorithm}[H]
  \caption{Pseudocódigo da Coletar e Processamento de dados brutos do \ac{gpr}.}
  \label{alg:pgpr}
  \SetAlgoLined
  \begin{algorithmic}[1]
    \STATE Coletar dados brutos do \ac*{gpr}
    \STATE Pré-processamento dos dados:
    \STATE Inicializar uma imagem vazia para armazenar a imagem otimizada.
    \FOR{$i=1$ até $i=4$ (i)}
    \STATE Inicializar uma imagem vazia para armazenar a imagem otimizada.
    \ENDFOR
    \STATE Pré-processamento:
    \STATE Realizar filtragem de imagem adicional.
    \STATE Aplicar técnicas de realce de recursos.
    \STATE Gerar visualizações finais da imagem otimizada.
    \STATE Atualize T se houver uma solução melhor
    \STATE Exibir a imagem otimizada para análise e interpretação.
    \STATE Salvar a imagem otimizada, se necessário.
    \STATE \Return{Imagem processada.}
  \end{algorithmic}
\end{algorithm}

%1. Coletar dados brutos do GPR.

%2. Pré-processamento dos dados:
%   a. Aplicar correções de ganho e calibração.
 %  b. Remover ruído indesejado.
  % c. Realizar correções de distorção.

%3. Inicializar uma imagem vazia para armazenar a imagem otimizada.

%4. Para cada ponto na grade da imagem:
 %  a. Definir uma região de interesse (ROI) em torno do ponto.
  % b. Calcular métricas de qualidade da imagem dentro da ROI, como contraste e nitidez.
   %c. Ajustar os parâmetros de processamento, como filtro e equalização, com base nas métricas de qualidade.
   %d. Aplicar os parâmetros de processamento otimizados à ROI.
   %e. Atualizar a imagem otimizada com os resultados do processamento.

%5. Repetir o passo 4 até que todos os pontos da imagem tenham sido processados.

%6. Pós-processamento opcional:
 ¨%  a. Realizar filtragem de imagem adicional.
 %  b. Aplicar técnicas de realce de recursos.
  % c. Gerar visualizações finais da imagem otimizada.

%7. Exibir a imagem otimizada para análise e interpretação.

%8. Salvar a imagem otimizada, se necessário.

%9. Fim do algoritmo.

Métodos avançados de processamento de dados do \ac*{gpr}, como filtragem espacial e temporal, silenciamento seletivo, filtragem de inclinação, deconvolução e análise de semelhança de velocidade, resultam em dados substancialmente diferentes das informações brutas. Além disso, operações específicas para o \ac*{gpr}, como remoção de fundo e mistura de antenas de múltiplas frequências, são aplicadas.

A evolução do \ac*{gpr} é documentada nos Anais de Conferências de \ac*{gpr} e fornece informações valiosas sobre o desenvolvimento do processamento de dados do \ac*{gpr}. A discussão aqui se concentra em questões essenciais do \ac*{gpr}, que diferem significativamente do processamento sísmico.

\subsection{Dewow}
\label{subsubsec:dewow}

Processo que é usado para suprimir um efeito chamado "wow" que ocorre nos dados do Radar de Penetração no Solo \ac{gpr} devido à proximidade do receptor com o transmissor. O "wow" é causado pela energia de baixa frequência nos campos próximos ao transmissor, que faz com que o nível base do sinal recebido se curve para cima ou para baixo. O "dewow" envolve a aplicação de um filtro temporal de alto amortecimento para eliminar esse efeito, e pode ser realizado em sistemas analógicos mais antigos ou digitalmente em sistemas modernos.

\subsection{Ganho de Tempo}
\label{subsubsec:ganhodetempo}

O "ganho de tempo", que é usado para equalizar as amplitudes dos sinais de radar de \ac*{gpr}, uma vez que esses sinais se atenuam rapidamente à medida que se propagam no solo. O ganho de tempo envolve a aplicação de uma função de ganho dependente do tempo para compensar essa rápida queda nos sinais de radar de maiores profundidades. A atenuação no solo pode variar significativamente, portanto, o ganho de tempo é aplicado para acomodar essas variações e permitir a exibição adequada dos dados de \ac*{gpr} versus tempo. O ganho de tempo pode ser aplicado de várias maneiras e deve ser escolhido com base em um modelo físico para minimizar artefatos no processo.

\subsection{Deconvolução}
\label{subsubsec:deconvolução}

Processo que visa maximizar a largura de banda e reduzir a dispersão de pulsos para aumentar a resolução em dados de \ac*{gpr}. No entanto, geralmente, a deconvolução de dados de \ac*{gpr} não oferece muitos benefícios significativos, pois o pulso padrão do \ac*{gpr} já é o mais curto e comprimido possível para as condições dadas de largura de banda e relação sinal-ruído. A deconvolução pode ser útil em casos de reverberação extrínseca ou reverberação do sistema.

Um tópico relacionado é a filtragem inversa Q, que tenta compensar a rápida atenuação das frequências mais altas em dados de \ac*{gpr} à medida que a profundidade aumenta. Isso é importante para manter a resolução em maiores profundidades. \citeonline{irving2003removal} fornecem uma discussão detalhada sobre a filtragem inversa Q.

\subsection{Migração}
\label{subsubsec:migração}


A migração é uma técnica de deconvolução espacial que visa remover a direcionalidade da fonte e do receptor dos dados de reflexão. Seu objetivo é reconstruir uma representação precisa e geométrica da distribuição de refletividade do subsolo, mas para isso, é necessário ter conhecimento da estrutura de velocidade. A migração é frequentemente um processo interativo, onde a velocidade de fundo é ajustada iterativamente para otimizar a imagem resultante conforme figura \ref{fig:fig06} . Existem várias formas de migração possíveis, e seu uso adequado no \ac*{gpr} visa compensar as características direcionais das antenas. Embora muitos princípios de migração sísmica se apliquem ao \ac*{gpr}, esse método também lida com o desafio único da topografia em relação à profundidade da exploração. A consideração da topografia é discutida por \citeonline{lehmann2000topographic}.

\begin{figure}[H]
  \centering
  \includegraphics[width=1\linewidth]{lib/fig/gpr/fig06.png}
  % Contém Fonte: 
  \caption{Imagens obtidas com um sistema de 50 MHz que atravessam dois túneis rodoviários, antes e após a migração.}
\source{\cite{jol2008ground}}
  \label{fig:fig06}
\end{figure}

\subsection{Correção topográfica}
\label{subsubsec:topografica}

Devido à pequena profundidade de exploração do \ac*{gpr}, a compensação da topografia frequentemente é importante. Para variações superficiais menores, o deslocamento temporal das trajetórias de dados pode compensar em grande parte as variações topográficas. O exemplo de perfil de reflexão na figura \ref{fig:fig10}  mostra a estratigrafia e uma reflexão da superfície da água. O perfil começa em uma área plana e sobe uma encosta. Os dados foram compensados por deslocamento temporal em relação à elevação ao longo da linha do perfil, o que torna a reflexão da superfície da água essencialmente plana.

\begin{figure}[H]
  \centering
  \includegraphics[width=1\linewidth]{lib/fig/gpr/fig10.png}
  % Contém Fonte: 
  \caption{Ilustração usando o deslocamento temporal dos traços para compensar variações topográficas. Observe que o nível freático é essencialmente plano na seção de radar de penetração no solo \ac{gpr}.}
\source{\cite{jol2008ground}}
  \label{fig:fig10}
\end{figure}